% ======================================================================
% Ontologie der Kontinua -- Core 1.3.3
% Cross-K-Modul: crossk_k8_k9.tex
% Cross-Level-Relationen zwischen K8 und K9
% ======================================================================

\subsubsection{K8–K9-Querverbindung}
\label{sec:crossk-k8-k9}

Der Übergang $K_8 \rightarrow K_9$ markiert die Entstehung des
\textbf{theoretisch-wissenschaftlichen und metaparadigmatischen Kontinuums}
aus dem zivilisatorisch-technologischen Kontinuum. Während $K_8$ populationsweite
Systeme, Infrastrukturen, symbolische Kodifizierung und technologische Zyklen
trägt, führt $K_9$ ein:
\begin{itemize}
  \item formale theoretische Rahmen,
  \item logische und methodische Achsen,
  \item Schwellen für Konsistenz, Evidenz und Erklärungskraft,
  \item Flüsse von Argumentation, Beweis und Kritik,
  \item Paradigmenzyklen und wissenschaftliche Entwicklung,
  \item rekursive Strukturen, die untere Stufen und sich selbst beschreiben können.
\end{itemize}

$K_9$ ist die erste Stufe, auf der Wissen zu einem expliziten strukturellen
Objekt mit eigenen Kontinuitätsbedingungen wird.

% ----------------------------------------------------------------------
\subsubsection{1. Stufen und ihre Rollen}

\paragraph{K8 (zivilisatorisch-technologisches Kontinuum).}
$K_8$ liefert:
\begin{itemize}
  \item symbolische Systeme, Schrift, Mathematik, formale Sprachen,
  \item Institutionen zur Wissensweitergabe und -akkumulation,
  \item Informationsnetze und kodierte Repositorien,
  \item rechnerische und technologische Fähigkeiten,
  \item Schwellen $\Theta_8$ (Energie, Logistik, informationelle Kohärenz).
\end{itemize}

Allerdings fehlt in $K_8$:
\begin{itemize}
  \item formal-meta-strukturelle Vergleiche von Theorien,
  \item logische Schwellen für Konsistenz oder Widerspruch,
  \item explizite Maße für erklärende oder prädiktive Stärke,
  \item strukturierte Zyklen wissenschaftlicher Evolution.
\end{itemize}

\paragraph{K9 (theoretisches Kontinuum).}
$K_9$ führt ein:
\begin{itemize}
  \item Achsen $A^9$ (Logik, Ontologie, Methodik, Interpretation),
  \item Potentiale $P^9$ (Erklärungskraft, Vorhersagegenauigkeit,
        Kohärenz, Sparsamkeit),
  \item Schwellen $\Theta^9$ (logische Konsistenz, empirische Angemessenheit,
        heuristische Tiefe, Gödel-Grenzen),
  \item Flüsse $J^9$ (Argumente, Beweise, Kritik, Paradigmenwechsel),
  \item Zyklen $C^9$ (Normale Wissenschaft, Anomalieakkumulation, Revolution),
  \item ein Maß $k(K_9)$ zur globalen theoretischen Kohärenz.
\end{itemize}

Somit ist $K_9$ eine meta-repräsentative Ebene oberhalb der zivilisatorischen
Achse.

% ----------------------------------------------------------------------
\subsubsection{2. Vererbte und neue Achsen sowie Schwellen}

\paragraph{Symbolische Vererbung.}
Symbolische Infrastrukturen von $K_8$ werden zum Substrat für theoretische
Unterscheidungen:
\[
A_{\mathrm{symbol}}(K_8) \rightarrow A^9.
\]
Schrift → Mathematik → Logik → formale Modelle.

\paragraph{Neue Achsen.}
$K_9$ führt ein:
\begin{itemize}
  \item logische Achsen (Gültigkeit, Schlussregeln),
  \item ontologische Achsen (Entitätstypen, strukturelle Annahmen),
  \item methodische Achsen (Modelle, Heuristiken, Paradigmen),
  \item Interpretationsachsen (konkurrierende Lesarten derselben Theorie).
\end{itemize}

\paragraph{Neue Schwellen.}
$\Theta^9$ umfasst:
\begin{itemize}
  \item \textbf{logische Konsistenzschwelle} — Verletzung führt zum Kollaps der Theorie,
  \item \textbf{empirische Schwelle} — Nichtkorrelation mit Evidenz,
  \item \textbf{heuristische Schwelle} — unzureichende Erklärungstiefe,
  \item \textbf{Paradigmen-Schwelle} — Inkompatibilität mit bestehenden Rahmen,
  \item \textbf{Gödel-Schwelle} — Unvollständigkeitsbeschränkungen.
\end{itemize}

Wenn Informationskohärenz $\Theta_I$ in $K_8$ versagt:
\[
\Theta_I^{(K_8)\mathrm{death}} \Rightarrow \Omega(K_9)=\varnothing.
\]

% ----------------------------------------------------------------------
\subsubsection{3. Cross-Level-Flüsse, Zyklen und Spannungen}

\paragraph{Flüsse.}
Theoretische Flüsse $J^9$ erweitern symbolische Flüsse von $K_8$:
\[
J^9 = (J_{\mathrm{argument}},\ J_{\mathrm{proof}},\
       J_{\mathrm{critique}},\ J_{\mathrm{reinterpret}}).
\]

Diese Flüsse verändern $A^9$, $P^9$ und die Grenzen $\partial\Omega(K_9)$.

\paragraph{Zyklen.}
Wissenschaftliche Zyklen $C^9$ umfassen:
\[
C_{\mathrm{science}}:\
\text{Theorie} \rightarrow \text{Vorhersage} \rightarrow \text{Experiment}
\rightarrow \text{Anomalie} \rightarrow \text{Revision} \rightarrow \text{Theorie}.
\]

Ein theoretisches Kontinuum existiert nur, wenn mindestens ein solcher stabiler
Zyklus besteht.

\paragraph{Spannung.}
Cross-Level-Spannung:
\[
T_{8,9} =
f(\Theta^9 - P^9,\ A^9 - A_{\mathrm{symbol}},\
  J^9 - J_{\mathrm{info}},\ \Theta_8 - \Theta^9).
\]
Überschreitet $T_{8,9}$ die dimensionsbezogene Schwelle, brechen Paradigmen in
nicht-theoretische symbolische Strukturen zurück.

% ----------------------------------------------------------------------
\subsubsection{4. Geburts- und Todesbedingungen über die Stufen}

\paragraph{Geburt von \texorpdfstring{$K_9$}{K_9}.}
Das theoretische Kontinuum entsteht, wenn:
\[
T_8 > \Theta_{8,\mathrm{dim}}
\quad\text{und}\quad
A^9 \in M_9 \setminus A(K_8),
\]
d.\,h. neue logische/ontologische Achsen auftreten, die nicht auf Technik oder
zivilisatorische Organisation zurückführbar sind.

Erweiterung:
\[
\Omega(K_9) = \Omega(K_8) \cup \Delta\Omega_{\mathrm{theory}}.
\]

\paragraph{Todesausbreitung.}
Wenn $\Theta_8$ versagt (zivilisatorischer Kollaps), verliert $K_9$ sein
Substrat.

Wenn $\Theta^9$ versagt (logischer Widerspruch, empirische Falsifikation,
Paradigmenbruch), kollabiert das theoretische Kontinuum, während $K_8$ bestehen bleibt:
\[
\Omega(K_9) \rightarrow \varnothing.
\]

% ----------------------------------------------------------------------
\subsubsection{5. Konzeptionelle Beispiele}

\paragraph{Formalierung von Wissen.}
Technisch-symbolische Systeme aus $K_8$ ermöglichen mathematische und logische
Strukturen, die Achsen in $K_9$ bilden.

\paragraph{Wissenschaftlicher Revolution-Zyklus.}
Wenn Anomalien über $\Theta^9_{\mathrm{heuristic}}$ akkumulieren:
\[
T_{8,9} > \Theta^9_{\mathrm{paradigm}}.
\]
tritt ein Paradigmenwechsel auf, der eine neue Region von $\Omega(K_9)$ schafft.

\paragraph{Kollapsfall.}
Sinkt die Informationskohärenz in $K_8$ unter $\Theta_I$, wird
theoretische Produktion unmöglich — charakteristisch für zivilisatorischen
Kollaps.

Damit kann der K8→K9-Übergang in komprimierter Form als Entstehung formaler
Theorien, Logik und wissenschaftlicher Evolution aus technologischen, symbolischen
und institutionellen Infrastrukturen zusammengefasst werden.
